\subsection{An offset must explain the individual rows}
\label{sec:apple-coupled-completion}

The required offset can be tested without pretending that its transaction has
been observed. Suppose a separate economic obligation would justify an
expense and an unpaid liability of $q$ in History 1 and $q-1{,}000$ in
History 2, with $q\geq1{,}000$. These are alternative amounts in two histories;
they are not a release booked within either history. For a numerical example,
take $q=2{,}000$: each branch debits expense and credits its own liability,
for 2,000 or 1,000 respectively. The difference supplies the 1,000 lower
liability and 1,000 higher profit required of History 2's remainder.
The economic obligation is an additional premise. Subtracting reported
totals has not established that it exists.

Even conditional on that premise, a matched total is insufficient. Assume
both the retirement loss and the new expense enter product cost, and both
the capital payable and the new liability enter the same non-tax current
liability category. History 2's extra loss then offsets its smaller expense
within that expense row, and its larger capital payable offsets its smaller
expense liability within that liability row. Other operations, distributions,
share counts and OCI remain fixed in this comparison. Equal tax expense and
tax balances are a further condition: equal pretax profit does not imply
equal deductions, tax bases or deferred tax.

The indirect cash-flow statement exposes another restriction. If the
expense liability is treated as an ordinary operating accrual, its 1,000
smaller increase reduces the working-liability adjustment by 1,000. The
extra retirement loss increases the noncash addback by 1,000. Total operating
cash flow is unchanged, but the separately reported Other adjustment and
Other current and non-current liabilities adjustment each change. That
construction fails to reproduce Apple's detailed cash-flow statement,
even though it balances and matches total operating cash
\citep[printed p.~33]{apple2024}.

A different conditional reporting map identifies the expense liability's
recognition as a separate noncash adjustment. Its amount must then be
excluded from the working-liability adjustment to avoid adding it back
twice. Let $\alpha$ be the fraction of this hypothetical recognition
allocated to the separate noncash adjustment; the remaining $1-\alpha$
enters the operating-liability adjustment. This parameter describes an
assumed allocation, not a probability or a freely selectable accounting
policy. With the retirement loss also assigned to Other, the differences
between History 2 and History 1 are
\[
\Delta\mathrm{Other}=1{,}000-\alpha\,1{,}000,
\qquad
\Delta\mathrm{WorkingLiabilities}=-(1-\alpha)\,1{,}000,
\qquad
\Delta\mathrm{CFO}=0.
\]
For $\alpha=0$, the two rows differ by $+1{,}000$ and $-1{,}000$.
For $\alpha=\tfrac12$, they still differ by $+500$ and $-500$.
Both row differences vanish only at $\alpha=1$ within this construction.
Including the full liability movement as well as the separate addback
would instead understate the History 2 minus History 1 operating-cash
difference by 1,000.

Under the common expense and liability mappings, equal taxes and the
disjoint noncash mapping with $\alpha=1$, all five annual statements are
unchanged by the combined transformation. Fixed shares also preserve EPS;
fixed OCI and distributions carry the profit equality into comprehensive
income and the equity statement. This is a conditional invariance result.
It assumes a feasible full-company History 1 exists and that the transformed
events remain admissible. It neither constructs that base nor proves that
either complete history is compatible with all notes. Assigning the two
expenses to different rows, the liabilities to different categories, or
different tax consequences immediately breaks the relevant equality.

The notes make these remaining conditions specific. Note 5 separates gross
PP\&E by asset class: the extra acquisition and retirement must cancel
within each affected class to leave that table unchanged. Note 6 separates
26,601 of current income taxes payable from 51,703 of other current
liabilities, and 9,254 of non-current income taxes payable from 36,634 of
other non-current liabilities. A total-liability offset cannot choose its
own location inside those amounts. Note 8 identifies finance-lease
right-of-use assets of 1,069 within PP\&E and distinguishes lease liabilities
and noncash lease acquisitions. Those disclosures do not identify the
hypothetical capital purchases or expense liability. Nor do Note 12's
purchase commitments and contingency assessment establish the assumed
obligation or its amount
\citep[printed pp.~39, 42, 45--46]{apple2024}.

The distinction therefore survives the repair: an algebraic offset is
available; a row-preserving transformation is available under stated
economic, tax and reporting assumptions; an actual Apple completion is
unestablished. To admit one, the historical model needs the underlying
obligation, gross asset classes, expense and liability assignments, tax
bases and noncash treatment, together with a feasible rest of the company.
The subsequent quarter additionally needs settlement dates and surviving
asset cohorts. Identical annual statements do not make those future
responses identical.
